\section{Behavior Under Algorithmic Redistricting}
\label{sec:behavior}

\subsection{Why Compactness Reduces EG}

The Efficiency Gap is elevated when one party's votes are systematically more
efficient than the other's---when one party wins many districts narrowly while
the other wins a few districts by large margins. This vote-distribution pattern
is the geometric signature of packing and cracking: cracking splits the target
party's voters across many just-losing districts (all votes wasted as losing-side
votes), while packing concentrates the remainder in a few landslide districts
(excess votes above 50\%+1 wasted as surplus).

Both cracking and packing require drawing district boundaries that reach across
natural geographic discontinuities. A cracked district must aggregate enough
opposing-party voters from dispersed locations to stay just below 50%; a packed
district must concentrate enough co-party voters from surrounding territory to
push the margin above 70--80\%. Both operations produce geographically irregular
boundaries---elongated districts, districts with narrow necks connecting distant
population clusters, or districts that wrap around geographic features to include
or exclude specific communities.

Compactness optimization in \textsc{bisect} penalizes exactly these boundary
configurations. By maximizing the ratio of area enclosed to perimeter (Polsby-Popper),
or equivalently by minimizing the perimeter given a fixed area, compactness
constraints prevent the geographic reaching that packing and cracking require.
Districts under \textsc{bisect} draw their populations from geographically
contiguous, approximately circular or convex regions. Because American residential
geography is heterogeneous at fine scales---partisan communities are not uniformly
distributed, but they are also not perfectly segregated---compactness constraints
produce districts with roughly competitive partisan compositions near the geographic
center of the state's partisan distribution.

The formal statement is: compactness optimization is sufficient to bound $|\text{EG}|$
near the geographic sorting baseline, because the only mechanism for creating large
EG (systematic packing and cracking) requires geographic configurations that
compactness penalizes. This is not a design guarantee---the algorithm does not
target any EG level---but an emergent property of the compactness constraint.

\subsection{Residual EG in Bisect Maps}

\textsc{Bisect} maps for NC ($k=14$) produce $|\text{EG}| < 0.03$, not $\text{EG}=0$.$^\dagger$
The residual EG reflects two sources:

\textbf{Geographic sorting.} Democratic voters are concentrated in the Triangle
(Research Triangle: Durham, Chapel Hill, Raleigh) and Charlotte urban cores.
Compactness constraints prevent drawing boundaries around these communities
to dilute them, which means compact districts centered on these cores will be
heavily Democratic. Any map with $k=14$ districts that includes compact districts
centered on these cores will have some positive EG, because the cores produce
large Democratic surpluses in a small number of districts.

\textbf{Arithmetic of bisection.} The \textsc{bisect} algorithm divides the state
into districts of approximately equal population by recursive geographic bisection.
When one party's voters are geographically concentrated, the bisection tends to
create one or two districts with strong partisan skew (the concentrated area)
and the remainder with weaker skew. The resulting vote-share distribution
produces a small but non-zero EG.

The key finding is that the residual EG in \textsc{bisect} maps ($\approx 0.02$)
is substantially below the enacted map EG ($\approx 0.09$) and reflects geography
rather than deliberate manipulation. The gap between the two values ($\approx 0.07$)
is the mapmaker contribution---the additional EG beyond what a neutral process
produces for the same state.

\footnotetext{$^\dagger$ Single-run \textsc{bisect} result using VEST 2020
presidential precinct returns, standard-bisect structure, geographic weights.}
